\section{1.3 실습 환경 세팅: Gymnasium}
\begin{frame}{왜 표준 인터페이스가 중요한가}
  이번 학기 실습은 \kb{Gymnasium}(Towers et al., 2024) ---
  옛 OpenAI Gym의 후신 --- 위에서 진행.
  \begin{itemize}
    \item 어떤 환경이든 ``현재 상태를 관측(\texttt{reset}),
      행동을 하나 취하면(\texttt{step}) \textbf{다음 상태$\cdot$보상$\cdot$
      종료 여부}가 돌아온다''는 \textbf{같은 인터페이스}.
    \item 2016년 Gym(Brockman et al., 2016) 이전: 각 연구자가 각자
      시뮬레이터를 각자 감쌌고, 재현하려면 ``이 사람의 환경 API를
      읽는 것''부터.
    \item Gym은 \textbf{``모든 환경이 동일한 두 함수
      (\texttt{reset}/\texttt{step})를 제공해야 한다''는 규칙 하나}로
      수천 개의 환경을 하나의 코드베이스 위에서 비교 가능하게.
    \item 이름 그대로 \textit{gym(체육관)} --- ``에이전트를 훈련시키러
      들어가는 곳''. 이후 \textbf{Gymnasium 포크}가 사실상 표준.
      이번 책의 코드는 \textbf{Gymnasium 1.x} 기준.
  \end{itemize}
\end{frame}
